\documentclass{article}
\usepackage[utf8]{inputenc}
\newcommand{\comment}[1]{ }
\usepackage{tabularx}
\usepackage{verbatim}
\usepackage{hyperref}
\usepackage{mathtools}
\usepackage{float}
\usepackage{subcaption}
\usepackage{textcomp}
\usepackage[a4paper, margin={1in}]{geometry}
\usepackage[english]{babel}
\usepackage{graphicx}
\usepackage{amsmath}
\usepackage{amsfonts}
\usepackage{fontspec}
\setmonofont{FreeMono}
\usepackage{minted}
\setminted{breaklines=true}
\newmintinline[lean]{lean4}{bgcolor=white}
\newminted[leancode]{lean4}{fontsize=\footnotesize}
\usemintedstyle{tango}

\frenchspacing
\usepackage{amsthm}
\usepackage{authblk}
\usepackage{biblatex}
\usepackage{mathpartir}
\usepackage{stmaryrd}
\usepackage{amssymb}
\newtheorem{lemma}{Lemma}
\newtheorem{theorem}{Theorem}
\addbibresource{./refs.bib}

\newcommand{\llb}{\llbracket}
\newcommand{\rrb}{\rrbracket}
\newcommand{\la}{\langle}
\newcommand{\ra}{\rangle}
\newcommand{\arr}{\rightarrow}
\newcommand{\sort}[1]{\text{Sort}\ #1}
\newcommand{\barr}{\arr_{\beta}}
\newcommand{\barrn}{\arr_{\beta}^{+}}
\newcommand{\earr}{\arr_{\eta}}
\newcommand{\affarr}{\rightarrowtail}
\newcommand{\affetaarr}{\rightarrowtail_{\eta}}
\newcommand{\etaarr}{\rightarrow_{\eta}}
\newcommand{\dom}[1]{\text{dom}(#1)}
\newcommand{\domtwo}[1]{\text{dom}_{2}(#1)}
\newcommand{\comb}[2]{\text{#1}(#2)}
\newcommand{\arity}[2]{\text{arity} #1 = #2}
\newcommand{\efn}{\text{EtaFn}}
\newcommand{\ddom}{\text{Domain}}
\newcommand{\canon}{\text{Canonicalize}}
\newcommand{\dddom}{\text{Domain}_{2}}

\newcolumntype{Y}{>{\centering\arraybackslash}X}

\title{DCC Beta Reduction Rules}
\author[1]{Lexi Aiello}

\begin{document}

\section{\(\eta\) Domain}

\begin{figure}[h!]
  \centering
  \resizebox{\textwidth}{!}{
    \begin{tabularx}{\textwidth}{l c r}
      \inferrule[Eql]
      {\dom{B} = A}
      {\dom{\comb{Eql}{B}} \coloneq !(A)} &
      \inferrule[Unbox]
      { }
      {\dom{*(A)} \coloneq !(A)}&
      \multicolumn{1}{r}{\inferrule[Eql1]
      {B x \barrn A}
      {\dom{\comb{Eql}{B} \&(x) \coloneq !(A)}}}\\\\
      \inferrule[Refl]
      {\dom{B} = A}
      {\dom{\comb{refl}{B}} \coloneq !(A)} &
      \inferrule[Refl1]
      {B x \barrn A}
      {\dom{\comb{refl}{B} \&(x)} \coloneq !(A)}&
      \multicolumn{1}{r}{\inferrule[Unmap]
      {\dom{x} = A}
      {\dom{*\la x \ra} \coloneq !(A)}}\\\\
      \inferrule[Fmap]
      {\dom{x} = A}
      {\dom{\$ \la x \ra} \coloneq !(A)}&
      \inferrule[Map]
      {\dom{x} = A}
      {\dom{! \la x \ra} \coloneq !(A)}&
      \multicolumn{1}{r}{\inferrule[Id]
      { }
      {\dom{\comb{Id}{A}} \coloneq A}}\\\\
      \inferrule[Fstc]
      { }
      {\dom{\comb{Fstc}{B}} \coloneq \comb{Coup}{B}}&
      \inferrule[Sndc]
      { }
      {\dom{\comb{Sndc}{B}} \coloneq \comb{Coup}{B}}&
      \multicolumn{1}{r}{\inferrule[Curr]
      { }
      {\dom{\Lambda (A, B, x)} \coloneq A}}\\\\
      \inferrule[Curr1]
      {B x \barrn B'}
      {\dom{\Lambda (A, B, x) y} \coloneq B'}&
      \inferrule[Quot]
      { }
      {\dom{`(A, x)} \coloneq A}&
      \multicolumn{1}{r}{\inferrule[Curr']
      { }
      {\dom{!\Lambda (A, B, x)} \coloneq !(A)}}\\\\
      \inferrule[Curr'1]
      {B x \barrn B'}
      {\dom{!\Lambda (A, B, x) y} \coloneq !(B')}&
      \inferrule[Pair]
      {\dom{x} = A}
      {\dom{\llb B \rrb \la x, y \ra} \coloneq A}&
      \multicolumn{1}{r}{\inferrule[Pi]
      {\dom{B} = A}
      {\dom{\Pi(B, C)} \coloneq A}}\\\\
      \inferrule[Eval]
      { }
      {\dom{\comb{Eval}{A, B, x}} \coloneq A}&
      \multicolumn{1}{c}{
      \inferrule[Eval1]
      {B y \barrn B'}
      {\dom{\comb{Eval}{A, B, x} y} \coloneq B'}}&
      \multicolumn{1}{r}{\inferrule[Comp']
      {\dom{g} = A}
      {\dom{f \circ' g} \coloneq A}}\\\\
      \multicolumn{1}{l}{
      \inferrule[Comp]
      {\dom{g} = !(A)}
      {\dom{f \circ g} \coloneq !(A)}}&
      \multicolumn{1}{c}{
      \inferrule[Comp2]
      {\dom{g} = !(A)}
      {\dom{f \circ\circ g} \coloneq !(A)}}&
      \multicolumn{1}{r}{ }
    \end{tabularx}
  }
  \caption{Where \(\earr\) represents \(\eta\)-reduction and \(\barrn\) represents \(\beta\)-normal reduction and \text{dom} syntactically retrieves the \(\beta\)-reduced domain of a combinator from environment. All domains are \(\beta\)-normal. \(\Sigma\) is omitted, since it has no \(\eta\)-reductino aside from reducing its points. We also model domains as a relation \(\ddom x A\) with completeness and soundness with respect to \(\dom{x}\).}
\end{figure}

\begin{figure}[H]
  \centering
  \resizebox{\textwidth}{!}{
    \begin{tabularx}{\textwidth}{l c r}
      \inferrule[Quot]
      { }
      {\domtwo{`(A, `(B, x))} \coloneq (A, B)}&
      \multicolumn{1}{c}{
      \inferrule[Id]
      {B x \barrn A}
      {\domtwo{\comb{Id}{\Pi(B, C) x}} \coloneq (\Pi(B, C) x, A)}}&
      \multicolumn{1}{r}{
      \inferrule[Comp']
      {\dom{g} = A, \domtwo{f} = (B, C)}
      {\domtwo{f \circ' g} \coloneq (A, C \circ' f)}}\\\\
      \inferrule[Eql]
      {\dom{B} = A}
      {\domtwo{\comb{Eql}{B}} \coloneq (A, B)}&
      \multicolumn{1}{c}{
      \inferrule[Refl]
      {\dom{B} = A}
      {\domtwo{\comb{refl}{B}} \coloneq (A, B)}}&
      \multicolumn{1}{r}{
      \inferrule[Curr]
      {\dom{x} = \comb{Coup}{B}}
      {\domtwo{\Lambda(A, B, x)} \coloneq (A, B)}}\\\\
      \multicolumn{3}{l}{
      \inferrule[Curr']
      {\dom{x} = !(\comb{Coup}{B})}
      {\domtwo{!\Lambda(A, B, x)} \coloneq (!(A), !\la B \ra)}}
    \end{tabularx}
  }
  \caption{Arity 2 function domains. The second projection of \(\dom{2}{x}\) is a function of the first term argument. Double quotation is \(\eta\)-reducible and considered to have arity 2 when double quoted. Similarly, when \(Id\) accepts a function, it has arity 2.}
\end{figure}

\section{\(\eta\) Holes}

\begin{figure}[H]
  \centering
  \resizebox{\textwidth}{!}{
    \begin{tabularx}{\textwidth}{l c r}
      \multicolumn{1}{r}{\inferrule[\(\eta\)-unary functions]
      {h_{arr} : \comb{arity}{x} = 1, h_{dom} : \ddom x A}
      {tm_{1} h_{arr} h_{dom} : \efn x [A]}}&
      \multicolumn{2}{l}{\inferrule[\(\eta\)-binary functions]
      {h_{arr} : \comb{arity}{x} = 2, h_{dom} : \dddom x A B}
      {tm_{2} h_{arr} h_{dom} : \efn x [A, B]}}\\\\
      \multicolumn{3}{c}{
      \inferrule[Constant holes]
      {h : (\neg \exists B, A = \comb{Coup}{B}) \land (\neg \exists B, A = !(B)), \text{level} : \mathbb{N}}
      {\text{pure} h : \text{Hole} A \text{level} (\text{const} (\text{Hole.mk} \text{level}))}}\\\\
      \multicolumn{3}{c}{\inferrule[Couple holes]
      {h_{dom} : \ddom B A, h_{1} : \text{Hole} A (\text{level} + 2) x, hr : (B eh_{1}) \barrn B', h_{2} : \text{Hole} B' (\text{level} + 1) y, \text{level} : \mathbb{N}}
      {\text{coup} h_{dom} h_{1} hr h_{2} : \text{Hole} \comb{Coup}{B} \text{level} \llb B \rrb (x, y)}}\\\\
      \multicolumn{3}{c}{\inferrule[Affine bang holes]
      {h_{1} : \text{Hole} A (\text{level} + 1) x, \text{level} : \mathbb{N}}
      {\text{bang} h_{1} : \text{Hole} !(A) \text{level} \&(x)}}
    \end{tabularx}}
\end{figure}

We also make use of relations \(\text{Domain} x [A]\) and \(\text{Domain}_{2} x [A, B]\) as proofs of function domains. Domains are computed by \(\text{dom} x\) and \(\text{dom}_2 x\) respectively. \(B\) is a function of the term argument \(y : A\).

\section{\(eta\) Canonicalization}

\begin{figure}[H]
  \centering
  \resizebox{\textwidth}{!}{
    \begin{tabularx}{\textwidth}{c}
      \inferrule[Id]
      {\efn x [A], (\neg \exists B C y, x = \Pi(B, C) y), \text{Hole} A y, x y \barrn y}
      {\canon e_f (\ddom.id (A \coloneq A))}\\\\
      \inferrule[Fstc]
      {\text{ddom} B A, \efn x [\comb{Coup}{B}], \text{Hole} \comb{Coup}{B} (\llb B \rrb (y, z)), x \llb B y \rrb (y, z) \barrn y}
      {\canon e_f (\ddom.fstc (B \coloneq B))}\\\\
      \inferrule[Sndc]
      {Domain B A, \efn x [\comb{Coup}{B}], \text{Hole} \comb{Coup}{B} (\llb B \rrb (y, z)), x \llb B y \rrb (y, z) \barrn z}
      {\canon e_f (\ddom.sndc (B \coloneq B))}\\\\
      \inferrule[Fstc']
      {Domain B A, \efn x [!(\comb{Coup}{B})], \text{Hole} !(\comb{Coup}{B}) (\&(\llb B \rrb (y, z))), x \&(\llb B y \rrb (y, z)) \barrn y}
      {\canon e_f (\ddom.unmap (\ddom.fstc (B \coloneq B)))}\\\\
      \inferrule[Sndc']
      {\dom B A, \efn x [!(\comb{Coup}{B})], \text{Hole} !(\comb{Coup}{B}) (\&(\llb B \rrb (y, z))), x \&(\llb B y \rrb (y, z)) \barrn z}
      {\canon e_f (\ddom.unmap (\ddom.sndc (B \coloneq B)))}\\\\
      \inferrule[Par]
      {Domain A y, ((\neg \exists A', A = \comb{Coup}{A'}) \land e_f (\neg \exists A', A = !(A'))), \efn x [A], \text{Hole} A y, x y \barrn x' y, \efn x' [A]}
      {\canon e_f e_f'}\\\\
      \inferrule[Eval]
      {Domain B A, \efn x [\comb{Coup}{B}], \text{Hole} \comb{Coup}{B} \llb B \rrb (y, z), x \llb B \rrb (y, z) \barrn x' y z, \efn x' [A, B]}
      {\canon e_f (\ddom.eval (A \coloneq A) (B \coloneq B) (x \coloneq x'))}\\\\
      \inferrule[Quot]
      {\efn x [A], \text{Hole} A y, x y \barrn e', \neg \exists c, (\text{const} (\text{hole} c)) \in e'}
      {\canon e_f (\ddom.quot (B \coloneq A) (x \coloneq e'))}\\\\
      \inferrule[$\text{Id}_{2}$]
      {\efn x [\Pi(f_A, f_B) \text{env}, `(\Pi(f_A, f_B) \text{env}, B)], \text{Hole} (\Pi(f_A, f_B) \text{env}) y, \text{Hole} B z, x y z \barrn y z}
      {\canon e_f (\ddom.id (A \coloneq \Pi(f_A, f_B) \text{env}))}\\\\
      \inferrule[$\text{Quot}_{2}$]
      {\efn x [A, `(A, B)], \text{Hole} A y, \text{Hole} B z, x y z \barrn e', (\neg \exists c, (\text{const} (\text{hole} c)) \in e')}
      {\canon e_f (\ddom.quot (A \coloneq A) (B \coloneq B) (x \coloneq e'))}\\\\
      \inferrule[Bang Quot]
      {\efn x [!(A), ! \la B \ra], \text{Hole} !(A) y, (e_h2 : \la B \ra y \barrn B'), \text{Hole} B' z, x y z \barrn e', (\neg \exists c, (\text{const} (\text{hole} c)) \in e')}
      {\canon e_f (\ddom.curr' (A \coloneq A) (B \coloneq B) (\ddom.quot (B \coloneq !(\comb{Coup}{B})) (x \coloneq e')))}
    \end{tabularx}}
  \caption{Hole-application observes function behavior. Functions are accordingly mapped to canonical operations which correspond to operations in the cartesian closed category. Holes are injected into the AST with a \text{const} constructor.}
\end{figure}

\begin{comment}

  Eta reduction digram:
  - Explain how to calculate holes
    - Domain is fine, since domain just calculates the shape. Eta-reduction inside the domain is not necessary, just beta-reduction.

\end{comment}

\end{document}

